Las remesas familiares y los medios de pago (M2) presentan una relación lineal positiva durante el período 2015–2025. A medida que el flujo de remesas aumenta, también se observa un crecimiento sostenido de los medios de pago, evidenciando una asociación entre ambas variables. Este comportamiento es consistente con el papel que desempeñan las remesas dentro de la economía guatemalteca, ya que una parte importante de estos recursos ingresa al sistema financiero y contribuye al incremento de la liquidez.
A partir de 2019 se aprecia una ligera dispersión en la nube de puntos, generando un patrón tipo abanico. Este comportamiento coincide con el período en el que las remesas familiares mostraron una mayor volatilidad, principalmente durante los años de la pandemia y la posterior recuperación económica. No obstante, la tendencia general continúa siendo creciente, por lo que la asociación positiva entre ambas variables se mantiene.
En términos de política económica, la relación observada sugiere que el seguimiento de las remesas familiares resulta relevante para evaluar la evolución de la liquidez del sistema financiero. Debido a que las remesas están expuestas a factores externos, como el desempeño económico de los países de origen de los migrantes y cambios en el contexto internacional, es recomendable que las autoridades monetarias y las instituciones financieras monitoreen su comportamiento como un indicador complementario para la toma de decisiones.